\documentclass[12pt]{exam}

\usepackage[utf8]{inputenc}
\usepackage[T1]{fontenc}
\usepackage{lmodern}
\usepackage[margin=1in]{geometry}
\usepackage{amsmath,amssymb}
\usepackage{array}
\usepackage{enumitem}
\usepackage{fancyvrb}

% ---- subpart labels:  a.  b.  c.  (instead of the default (a)) ----
\renewcommand{\thepartno}{\alph{partno}}
\renewcommand{\partlabel}{\thepartno.}
\renewcommand{\questionlabel}{\thequestion.}

% ---- running header (Name / SID) and footer, on every page ----
\pagestyle{headandfoot}
\header{Name \rule[-2pt]{6.5cm}{0.4pt}}{}{SID \rule[-2pt]{6cm}{0.4pt}}
\footer{Spring 2026}{\thepage}{CS/EE 147 Midterm}
\headrule

% ---- answer-box helpers ----
\newcommand{\Ans}{\underline{Answer:}\ }
\newcommand{\smallbox}{\framebox[5cm]{\rule{0pt}{1.1cm}}}
\newcommand{\tfbox}{\framebox[1.7cm]{\rule{0pt}{1.7cm}}}
\newcommand{\widebox}[1][9.5cm]{\framebox[#1]{\rule{0pt}{1.1cm}}}
\newcommand{\paragraphbox}[1]{\framebox[\linewidth]{\rule{0pt}{#1}}}
\newcommand{\mc}{$\bigcirc$\ }   % blank multiple-choice bubble

\setlength{\parindent}{0pt}

\begin{document}

% =====================================================================
%  TITLE PAGE
% =====================================================================
\begin{center}
  {\Large\bfseries CS/EE 147}\\[2pt]
  {\Large\bfseries GPU Architecture and Parallel Programming}\\[18pt]
  {\large\bfseries Spring 2024}\\[14pt]
  {\large\bfseries Midterm}\\[14pt]
  {\large\bfseries Time Allowed: 80 Minutes}
\end{center}

\vspace{10pt}

\begin{center}
\renewcommand{\arraystretch}{1.45}
\begin{tabular}{|>{\centering\arraybackslash}m{4cm}|>{\centering\arraybackslash}m{4cm}|>{\centering\arraybackslash}m{4cm}|}
\hline
\textbf{Question} & \textbf{Points Possible} & \textbf{Points Scored} \\ \hline
1            & 4  & \\ \hline
2            & 4  & \\ \hline
3            & 4  & \\ \hline
4            & 13 & \\ \hline
5            & 12 & \\ \hline
6--9 (MC)    & 12 & \\ \hline
10--17 (T/F) & 12 & \\ \hline
18           & 8  & \\ \hline
19           & 12 & \\ \hline
20           & 4  & \\ \hline
Total        & 85 & \\ \hline
\end{tabular}
\end{center}

\vspace{14pt}

\begin{itemize}[leftmargin=2.2em,itemsep=4pt]
  \item Closed book, closed internet
  \item Allowed 1 page letter-size cheatsheet, front and back
  \item Calculators allowed, but likely not necessary
\end{itemize}

\newpage

% =====================================================================
%  QUESTIONS
% =====================================================================
\begin{questions}

% ---------------- Q1 ----------------
\question \textbf{[4 points] CUDA Kernel Launch}\\
You want to launch a kernel function, \texttt{kernel()}, with 32 blocks of 256 threads per
block. Please write your answer in the provided box.
\begin{parts}
  \part {}[2 points] How many warps are in each block?\\[10pt]
        \Ans \smallbox
  \part {}[2 points] How many total threads are in your kernel?\\[10pt]
        \Ans \smallbox
\end{parts}

\vspace{8pt}
% ---------------- Q2 ----------------
\question \textbf{[4 points]} Assume that we want to extend the basic vector addition code so
that \textbf{each thread} computes \textbf{eight consecutive input elements}. (For example,
thread 0 computes index 0--7, thread 1 computes index 8--15, etc.) Assume that variable
\texttt{i} should be initialized with the index for the first element to be processed by a
thread. What should \texttt{i} be initialized to?

\begin{Verbatim}[xleftmargin=2em]
__global__ void vecAddKernel(float* A, float* B, float* C, int n) {
    int i = ????
    for(int j = 0; j < 8; j++)
        if( (i+j) < n ) C[i+j] = A[i+j] + B[i+j];
    ...
}
\end{Verbatim}

\vspace{4pt}
\texttt{int i = }\,\widebox

\vspace{8pt}
% ---------------- Q3 ----------------
\question \textbf{[4 points]} Can the above vector add program be made faster using shared
memory? Explain why or why not.\\[10pt]
\Ans \paragraphbox{2.6cm}

\vspace{8pt}
% ---------------- Q4 ----------------
\question \textbf{[13 points] Reduction}\\
For a Reduction kernel that operates on a total of 10000 elements using a block size of 256
threads, answer the following: (Assume each block calculates a \texttt{partialSum[]} as in
Assignment 2)
\begin{parts}
  \part {}[3 point] How many steps will the reduction kernel take?\\[10pt]
        \Ans \smallbox
  \part {}[3 points] For a na\"ive reduction implementation, how many steps are divergent?\\[10pt]
        \Ans \smallbox
  \part {}[3 points] For an optimized reduction implementation, how many steps are divergent?\\[10pt]
        \Ans \smallbox
  \part {}[4 points] For the kernel mentioned above, the warp distribution statistics look as
        follows --\\[4pt]
        W0: 0, W1: 20, W2: 20, W3: 0, W4: 20, W5: 0, W6: 0, W7: 0, W8: 20, W9: 0, W10: 0,
        W11: 0, W12: 0, W13: 0, W14: 0, W15: 0, W16: 20, W17: 0, W18: 0, W19: 0, W20: 0,
        W21: 0, W22: 0, W23: 0, W24: 0, W25: 0, W26: 0, W27: 0, W28: 0, W29: 0, W30: 0,
        W31: 0, W32: 300\\[4pt]
        Why are W1, W2, W4, W8, and W16 all equal to 20?\\[10pt]
        \Ans \paragraphbox{3.4cm}
\end{parts}

\newpage
% ---------------- Q5 ----------------
\question \textbf{[12 points] CUDA Streams.}\\
Consider the following code snippet of a CUDA Streams implementation of Vector Add.

\begin{Verbatim}[xleftmargin=1em]
for (int i=0; i<n; i+=SegSize*2) {
(A)  cudaMemcpyAsync(d_A0, h_A+i, SegSize*sizeof(float),..., stream0);
(B)  cudaMemcpyAsync(d_B0, h_B+i, SegSize*sizeof(float),..., stream0);
(C)  vecAdd<<<SegSize/256, 256, 0, stream0>>>(d_A0, d_B0,...);
(D)  cudaMemcpyAsync(h_C+i, d_C0, SegSize*sizeof(float),..., stream0);

(E)  cudaMemcpyAsync(d_A1, h_A+i+SegSize, SegSize*sizeof(float),..., stream1);
(F)  cudaMemcpyAsync(d_B1, h_B+i+SegSize, SegSize*sizeof(float),..., stream1);
(G)  vecAdd<<<SegSize/256, 256, 0, stream1>>>(d_A1, d_B1,...);
(H)  cudaMemcpyAsync(d_C1, h_C+i+SegSize, SegSize*sizeof(float),..., stream1);
}
\end{Verbatim}

\begin{parts}
  \part {}[4 points] Which two instructions would experience overlapped execution? Fill in the
        boxes below \textbf{with the letter} associated with the two instructions that overlap.\\[10pt]
        \underline{Answer 1:}\ \framebox[4cm]{\rule{0pt}{1.1cm}}\hfill
        \underline{Answer 2:}\ \framebox[4cm]{\rule{0pt}{1.1cm}}
  \part {}[5 points] If you are allowed to \textbf{move only one instruction} up or down in this
        order to maximize pipelining in this case, which instruction would be moved below which
        other instruction?\\[10pt]
        \underline{Instr:}\ \framebox[4cm]{\rule{0pt}{1.1cm}}\hfill
        \underline{Below:}\ \framebox[4cm]{\rule{0pt}{1.1cm}}
  \part {}[3 points] If \textbf{SegSize = 256} and n=8192, how many vecAdd kernels will be
        launched during this program?\\[10pt]
        \Ans \smallbox
\end{parts}

\newpage
% ---------------- MC: Q6-Q9 ----------------
\textbf{Multiple Choice. \textnormal{Please completely fill in the circle} $\bigcirc$
\textnormal{for your answer.} [12 points total]}

\vspace{6pt}
% ---------------- Q6 ----------------
\question {}[3 points] If a CUDA device's SM (streaming multiprocessor) can take up to 2048
threads and up to 32 thread blocks. Which of the following 2D block configurations would
result in the greatest number of threads in each SM with the minimum number of blocks?
\begin{itemize}[leftmargin=2.5em,itemsep=4pt,label={}]
  \item \mc 64x64
  \item \mc 32x32
  \item \mc 16x16
  \item \mc 8x8
\end{itemize}

% ---------------- Q7 ----------------
\question {}[3 points] How many streams do you need to fully utilize all of the kernel and copy
engines on the GPU?
\begin{itemize}[leftmargin=2.5em,itemsep=4pt,label={}]
  \item \mc 1
  \item \mc 2
  \item \mc 3
  \item \mc 4
\end{itemize}

% ---------------- Q8 ----------------
\question {}[3 points] As data is transferred from host to device using \texttt{cudaMemCopy} it
touches various types of hardware components (memories, buses, on-chip components, etc.). Order
the following hardware components in the order of being used during a \texttt{cudaMemCopy} H2D
operation.
\begin{itemize}[leftmargin=2.5em,itemsep=4pt,label={}]
  \item \mc Pinned Memory -> Host Memory -> DMA Engine -> Global Memory
  \item \mc Host Memory -> DMA Engine -> Pinned Memory -> Global Memory
  \item \mc Host Memory -> Pinned Memory -> DMA Engine -> Global Memory
  \item \mc Pinned Memory -> DMA Engine -> Host Memory -> Global Memory
  \item \mc None of the above
\end{itemize}

% ---------------- Q9 ----------------
\question {}[3 points] Which of the following memory types has the fastest access speed?
\begin{itemize}[leftmargin=2.5em,itemsep=4pt,label={}]
  \item \mc L2 Cache
  \item \mc L1 Cache
  \item \mc Register
  \item \mc Global Memory
  \item \mc Shared Memory
\end{itemize}

\newpage
% ---------------- T/F: Q10-Q15 ----------------
\textbf{True or False. \textnormal{For the following questions, fill the answer box with (T)
True or (F) False.} [16 points total]}

\vspace{10pt}
\renewcommand{\arraystretch}{1.3}
\noindent
\begin{tabular}{@{}p{0.40\linewidth}c@{\hspace{0.5cm}}p{0.40\linewidth}c@{}}
10.~\texttt{\_\_device\_\_} functions are only callable from the host & \tfbox
 & 13.~A SIMT core has four independent schedulers & \tfbox \\[1.4cm]
11.~cudamemcpy is a synchronous process. & \tfbox
 & 14.~All cudamemcpy (host to device) copies from pinned memory & \tfbox \\[1.4cm]
12.~All blocks in a grid can be executed independent to each other in any order & \tfbox
 & 15.~The maximum number of threads and blocks that can reside in an SM is fixed for each GPU vendor. & \tfbox \\[1.0cm]
\end{tabular}

\vspace{16pt}
\setcounter{question}{15}
% ---------------- Q16 ----------------
\question {}[6 points] You need to write a 2D kernel that operates on an image of size 500x800
(height (y-axis) x width(x-axis)). Assume your thread block size is 16x16.
\begin{parts}
  \part {}[2 points] What is the x-dimension of your grid size?\\[10pt]
        \Ans \smallbox
  \part {}[2 points] What is the y-dimension of your grid size?\\[10pt]
        \Ans \smallbox
  \part {}[4 point] How many warps will be divergent in this scenario?\\[10pt]
        \Ans \smallbox
\end{parts}

\newpage
% ---------------- Q17 ----------------
\question {}[12 points] Active mask. Assume a warp size of 8, and the following pseudo-code will
be executed. Fill in the active mask at each of the following basic block locations. Hint: the
active mask for a warp size of 8 when all threads are active will be 11111111.

\begin{Verbatim}[xleftmargin=2em]
int N = 6;
__global__ function()
{
 int index = threadIdx.x; //PC = A

 if(index < N)
 {
   //PC = B
   if(index != 0)
   { //PC = C
     . . .
   }
   else
   { //PC = D
     . . .
   }
   //PC = E
 }
 else
 { //PC = F
   . . .
 }
 __syncthreads(); //PC = G
}
\end{Verbatim}

\vspace{8pt}
\begin{tabular}{@{}l@{\hspace{1cm}}l@{}}
\underline{Answer} (PC = B):\ \framebox[4.2cm]{\rule{0pt}{0.85cm}} &
\underline{Answer} (PC = C):\ \framebox[4.2cm]{\rule{0pt}{0.85cm}} \\[12pt]
\underline{Answer} (PC = D):\ \framebox[4.2cm]{\rule{0pt}{0.85cm}} &
\underline{Answer} (PC = F):\ \framebox[4.2cm]{\rule{0pt}{0.85cm}} \\
\end{tabular}

\vspace{16pt}
% ---------------- Q18 ----------------
\question {}[4 points] Assume that a kernel is launched with 100 thread blocks, each with 512
threads. If a variable, s\_var, is declared as a shared memory variable, how many versions of
s\_var will be created throughout the lifetime of the execution of the kernel?\\[10pt]
\Ans \smallbox

\end{questions}

\end{document}
